% Options for packages loaded elsewhere
\PassOptionsToPackage{unicode}{hyperref}
\PassOptionsToPackage{hyphens}{url}
\PassOptionsToPackage{dvipsnames,svgnames,x11names}{xcolor}
\documentclass[
  11pt,
]{article}
\usepackage{xcolor}
\usepackage[a4paper,margin=2.4cm]{geometry}
\usepackage{amsmath,amssymb}
\setcounter{secnumdepth}{-\maxdimen} % remove section numbering
\usepackage{iftex}
\ifPDFTeX
  \usepackage[T1]{fontenc}
  \usepackage[utf8]{inputenc}
  \usepackage{textcomp} % provide euro and other symbols
\else % if luatex or xetex
  \usepackage{unicode-math} % this also loads fontspec
  \defaultfontfeatures{Scale=MatchLowercase}
  \defaultfontfeatures[\rmfamily]{Ligatures=TeX,Scale=1}
\fi
\usepackage{lmodern}
\ifPDFTeX\else
  % xetex/luatex font selection
\fi
% Use upquote if available, for straight quotes in verbatim environments
\IfFileExists{upquote.sty}{\usepackage{upquote}}{}
\IfFileExists{microtype.sty}{% use microtype if available
  \usepackage[]{microtype}
  \UseMicrotypeSet[protrusion]{basicmath} % disable protrusion for tt fonts
}{}
\makeatletter
\@ifundefined{KOMAClassName}{% if non-KOMA class
  \IfFileExists{parskip.sty}{%
    \usepackage{parskip}
  }{% else
    \setlength{\parindent}{0pt}
    \setlength{\parskip}{6pt plus 2pt minus 1pt}}
}{% if KOMA class
  \KOMAoptions{parskip=half}}
\makeatother
\usepackage{longtable,booktabs,array}
\newcounter{none} % for unnumbered tables
\usepackage{calc} % for calculating minipage widths
% Correct order of tables after \paragraph or \subparagraph
\usepackage{etoolbox}
\makeatletter
\patchcmd\longtable{\par}{\if@noskipsec\mbox{}\fi\par}{}{}
\makeatother
% Allow footnotes in longtable head/foot
\IfFileExists{footnotehyper.sty}{\usepackage{footnotehyper}}{\usepackage{footnote}}
\makesavenoteenv{longtable}
\setlength{\emergencystretch}{3em} % prevent overfull lines
\providecommand{\tightlist}{%
  \setlength{\itemsep}{0pt}\setlength{\parskip}{0pt}}
% ─────────────────────────────────────────────────────────────────────────
% latex-header.tex — academic house style for the project's design docs.
% Injected into a standalone .tex by md2tex.sh (pandoc --include-in-header -s).
% Target: Overleaf with the default pdfLaTeX engine (no system fonts, no
% fontspec). Classic scholarly look: Palatino (newpx) body, navy headings,
% booktabs tables (pandoc emits these by default), microtype, A4.
% ─────────────────────────────────────────────────────────────────────────

% Fonts — Palatino-style serif via newpx (bundled with TeX Live / Overleaf).
\usepackage[T1]{fontenc}
\usepackage{newpxtext}
\usepackage{newpxmath}
\usepackage[scaled=0.92]{helvet}   % Helvetica-like sans
\usepackage{inconsolata}           % clean monospace

\usepackage{microtype}             % refined spacing / justification
\usepackage{xcolor}
\definecolor{headingnavy}{HTML}{1F3864}
\definecolor{rulegray}{HTML}{2F2F2F}

% Navy, serif section headings (matches the Word template).
\usepackage{titlesec}
\titleformat{\section}{\Large\bfseries\color{headingnavy}}{\thesection}{0.6em}{}
\titleformat{\subsection}{\large\bfseries\color{headingnavy}}{\thesubsection}{0.6em}{}
\titleformat{\subsubsection}{\normalsize\bfseries\color{headingnavy}}{\thesubsubsection}{0.6em}{}
\titlespacing*{\section}{0pt}{1.4ex plus 0.4ex}{0.7ex}
\titlespacing*{\subsection}{0pt}{1.1ex plus 0.3ex}{0.5ex}

% Booktabs rules in the same gray as the Word version.
\usepackage{booktabs}
\usepackage{colortbl}              % enables \arrayrulecolor
\arrayrulecolor{rulegray}
\setlength{\heavyrulewidth}{1.1pt}
\setlength{\lightrulewidth}{0.5pt}
\renewcommand{\arraystretch}{1.25}
\usepackage{ragged2e}

% Unicode glyph that pdfLaTeX/newpx doesn't render on its own.
% (≤ and ≥ are already provided, so we don't redefine them here.)
\usepackage{newunicodechar}
\newunicodechar{→}{\ensuremath{\rightarrow}}

% NOTE: link colours are set via pandoc -V variables (colorlinks/linkcolor/
% urlcolor) — NOT \hypersetup here, because pandoc loads hyperref *after*
% this header, so calling \hypersetup here would be an undefined command.

% Report-style spacing: blank line between paragraphs, no indent.
\usepackage{parskip}
\usepackage{bookmark}
\IfFileExists{xurl.sty}{\usepackage{xurl}}{} % add URL line breaks if available
\urlstyle{same}
\hypersetup{
  colorlinks=true,
  linkcolor={headingnavy},
  filecolor={Maroon},
  citecolor={Blue},
  urlcolor={headingnavy},
  pdfcreator={LaTeX via pandoc}}

\author{}
\date{}

\begin{document}

\section{Requirements Specification}\label{requirements-specification}

\textbf{Project:} RAG-Based Retrieval for Archival Systems\\
\textbf{Version:} 1.1\\
\textbf{Date:} 2026-06-09\\
\textbf{Status:} Draft (Revised)

\subsection{1. Introduction}\label{introduction}

\subsubsection{1.1 Purpose}\label{purpose}

This document specifies the functional requirements for a
Retrieval-Augmented Generation (RAG) system that provides
natural-language access to institutional archival document collections.
Users query a conversational interface; the system retrieves relevant
document chunks and returns a grounded, cited response generated by a
locally hosted LLM.

The system is built as a \textbf{microservices architecture}, with
pipeline flows orchestrated by \textbf{Node-RED} and \textbf{Islandora}
(via Drupal) as the archival repository.

\subsubsection{1.2 Scope}\label{scope}

\textbf{In scope:}

\begin{itemize}
\tightlist
\item
  Document ingestion (text extraction, chunking, embedding, indexing via
  Islandora)
\item
  Natural language query processing and hybrid retrieval
\item
  LLM-based response generation with source attribution
\item
  Chat-based front-end (web extension and mobile)
\end{itemize}

\textbf{Out of scope:}

\begin{itemize}
\tightlist
\item
  OCR processing of scanned/image-based documents
\item
  LLM fine-tuning
\item
  Handwritten document processing
\item
  Cloud inference APIs (all inference is on-premises)
\end{itemize}

\subsubsection{1.3 Glossary}\label{glossary}

{\def\LTcaptype{none} % do not increment counter
\begin{longtable}[]{@{}
  >{\raggedright\arraybackslash}p{(\linewidth - 2\tabcolsep) * \real{0.5000}}
  >{\raggedright\arraybackslash}p{(\linewidth - 2\tabcolsep) * \real{0.5000}}@{}}
\toprule\noalign{}
\begin{minipage}[b]{\linewidth}\raggedright
Term
\end{minipage} & \begin{minipage}[b]{\linewidth}\raggedright
Definition
\end{minipage} \\
\midrule\noalign{}
\endhead
\bottomrule\noalign{}
\endlastfoot
RAG & Retrieval-Augmented Generation: combines LLM generation with
external document retrieval \\
Chunk & A semantically bounded segment of a source document stored as a
retrievable unit \\
Embedding & A dense numerical vector representing the semantic meaning
of a text chunk \\
Dense Retrieval & Vector similarity search using embeddings \\
Sparse Retrieval & Keyword-based search (e.g., BM25) \\
Hybrid Retrieval & Combination of dense and sparse retrieval
strategies \\
Reranker & A model that re-scores retrieved candidates by relevance to
the query \\
Source Attribution & Tracing a generated response back to the specific
source document(s) it was grounded in \\
Contextual Chunk & A chunk enriched with structural metadata (section
heading, document hierarchy, date) \\
Islandora & Open-source digital repository framework (via Drupal) used
as the archival backend \\
Node-RED & Flow-based orchestration tool used to wire the microservices
into ingestion and query pipelines \\
Microservice & An independently deployable service implementing one
stage of the RAG pipeline \\
Query Normalisation & Preprocessing step that standardises query text by
lowercasing, removing excess punctuation, and applying canonical
forms \\
Keyword Expansion & Augmenting a query with synonyms or related terms to
improve retrieval recall \\
Ingestion Pipeline (RAG) & The automated project-specific workflow that extracts,
chunks, embeds, and indexes documents from Islandora into the RAG
vector database. Distinct from Islandora's own internal ingestion
pipelines. \\
\end{longtable}
}

\subsection{2. Functional Requirements}\label{functional-requirements}

\subsubsection{FR-01: Natural Language Query
Interface}\label{fr-01-natural-language-query-interface}

\textbf{Description:} The system shall accept natural language queries
from users through a chat interface (web and mobile). Queries may be
phrased as questions, keywords, or conversational prompts.

\textbf{Priority:} High\\
\textbf{Testability:} Submit a natural language question; verify the
system returns a textual response within acceptable latency.

\subsubsection{FR-02: Query
Preprocessing}\label{fr-02-query-preprocessing}

\textbf{Description:} The system shall preprocess each incoming query
before retrieval by applying language detection, query normalisation,
and keyword expansion.

\textbf{Priority:} Medium\\
\textbf{Testability:} Submit queries in multiple languages and verify
correct language detection; verify normalised forms are used in
retrieval.

\subsubsection{FR-03: Embedding Generation for
Queries}\label{fr-03-embedding-generation-for-queries}

\textbf{Description:} The system shall generate a dense vector embedding
for each preprocessed query using a multilingual sentence transformer
model.

\textbf{Priority:} High\\
\textbf{Testability:} Verify embedding dimensions match the vector
database schema; semantically similar queries must produce geometrically
close vectors.

\subsubsection{FR-04: Hybrid Retrieval}\label{fr-04-hybrid-retrieval}

\textbf{Description:} The system shall retrieve candidate chunks using a
combination of:

\begin{itemize}
\tightlist
\item
  \textbf{Dense vector search}: semantic similarity via embeddings
\item
  \textbf{Sparse keyword search}: BM25 or equivalent
\item
  \textbf{Metadata / filter search}: document type, date range,
  collection tag
\end{itemize}

\textbf{Priority:} High\\
\textbf{Testability:} Run identical queries with each strategy
enabled/disabled; verify recall improves with hybrid mode active.

\subsubsection{FR-05: Re-ranking of Retrieved
Candidates}\label{fr-05-re-ranking-of-retrieved-candidates}

\textbf{Description:} The system shall apply a re-ranking model to the
initial set of retrieved chunks to improve precision before context
assembly.

\textbf{Priority:} Medium\\
\textbf{Testability:} Compare ranked vs.~unranked output; top-k results
after reranking must be more relevant than naive retrieval order.

\subsubsection{FR-06: Context Assembly}\label{fr-06-context-assembly}

\textbf{Description:} The system shall select and assemble the
top-ranked chunks into a coherent context window for prompt
construction, respecting the LLM's token limit.

\textbf{Priority:} High\\
\textbf{Testability:} Verify total assembled token count does not exceed
the configured context budget.

\subsubsection{FR-07: LLM Response
Generation}\label{fr-07-llm-response-generation}

\textbf{Description:} The system shall pass the assembled context and
user query to a locally hosted open-weight LLM and return the generated
response to the user.

\textbf{Priority:} High\\
\textbf{Testability:} Verify responses are generated without external
API calls; verify the response is grounded in the provided context.

\subsubsection{FR-08: Source
Attribution}\label{fr-08-source-attribution}

\textbf{Description:} Every generated response shall include citations
identifying the source document(s), section, and page (where available)
from which the supporting evidence was retrieved.

\textbf{Priority:} High\\
\textbf{Testability:} For any factual claim in the response, the cited
document must contain text that supports that claim.

\subsubsection{FR-09: Structured JSON Chunk
Storage}\label{fr-09-structured-json-chunk-storage}

\textbf{Description:} The ingestion pipeline shall store document chunks
as structured JSON objects preserving the document's semantic hierarchy
with at minimum the following fields:

\begin{itemize}
\tightlist
\item
  \texttt{context\_heading}: the section/chapter title the chunk belongs
  to
\item
  \texttt{context\_content}: the chunk's actual text
\item
  \texttt{context\_links}: references or related source identifiers
\item
  \texttt{metadata}: document ID, date, language, collection tag
\end{itemize}

\textbf{Priority:} High\\
\textbf{Testability:} Inspect stored chunks in the vector DB; all
required fields must be populated.

\subsubsection{FR-10: Document Ingestion
Pipeline}\label{fr-10-document-ingestion-pipeline}

\textbf{Description:} The system shall support ingestion of new
documents from the Islandora archival repository through an automated
pipeline that:

\begin{enumerate}
\def\labelenumi{\arabic{enumi}.}
\tightlist
\item
  Extracts text from structured digital documents (PDF, plain text,
  HTML)
\item
  Detects document language
\item
  Splits the document into contextual chunks using structure-aware
  chunking
\item
  Generates embeddings for each chunk
\item
  Stores chunks with metadata in a vector database
\end{enumerate}

\textbf{Note:} OCR processing of scanned image documents is out of
scope.

\textbf{Priority:} High\\
\textbf{Testability:} Ingest a known PDF; verify all chunks appear in
the vector database with correct fields and no mid-sentence chunk
boundaries.

\subsubsection{FR-11: Islandora Repository
Integration}\label{fr-11-islandora-repository-integration}

\textbf{Description:} The system shall integrate with Islandora (via
Drupal) as the primary archival document store.

\textbf{Priority:} High\\
\textbf{Testability:} Add a document to Islandora; verify the document
is ingested and appears in the vector database within the defined SLA.

\subsubsection{FR-12: Microservices
Architecture}\label{fr-12-microservices-architecture}

\textbf{Description:} Each stage of the RAG pipeline (text extraction,
chunking, embedding, retrieval, reranking, LLM generation) shall be
implemented as an independently deployable microservice.

\textbf{Priority:} High\\
\textbf{Testability:} Stop or replace any single service (e.g., swap the
embedding service) and verify the rest of the pipeline continues to
operate correctly.

\subsubsection{FR-13: Node-RED Pipeline
Orchestration}\label{fr-13-node-red-pipeline-orchestration}

\textbf{Description:} The system shall use Node-RED to orchestrate the
microservices into two flows:

\begin{itemize}
\tightlist
\item
  \textbf{Ingestion flow}: Islandora trigger → text extraction →
  language detection → chunking → embedding → indexing
\item
  \textbf{Query flow}: query preprocessing → hybrid retrieval →
  reranking → context assembly → LLM generation → source attribution
\end{itemize}

\textbf{Priority:} High\\
\textbf{Testability:} Inspect the Node-RED flow; verify each pipeline
stage maps to a node, and that end-to-end flow executes correctly for
both ingestion and query paths.

\subsubsection{FR-14: Multilingual
Support}\label{fr-14-multilingual-support}

\textbf{Description:} The system shall support document retrieval and
response generation across multiple languages.

\textbf{Priority:} Medium\\
\textbf{Testability:} Ingest documents in German and English; submit
queries in both languages; verify relevant results are returned for
each.

\subsubsection{FR-15: Session Context / Follow-Up
Questions}\label{fr-15-session-context-follow-up-questions}

\textbf{Description:} The system shall maintain conversational context
within a session so that follow-up questions can reference prior turns
without the user restating context.

\textbf{Priority:} Medium\\
\textbf{Testability:} Ask a follow-up question referencing a prior
answer in the same session; verify the system resolves the reference
correctly.

\subsubsection{FR-16: Drupal Front-End (AI
Chatblock)}\label{fr-16-drupal-front-end-ai-chatblock}

\textbf{Description:} The system shall provide a chat-based front-end
within Drupal (the same platform that hosts Islandora) using the AI
Chatblock module (or its successor \texttt{ai\_assistant\_ui}).

\textbf{Priority:} Medium\\
\textbf{Testability:} Submit a query from the Drupal front-end; verify
the response is rendered correctly and citations are displayed.

\subsubsection{FR-20: Streaming Query API
(OpenAI-Compatible)}\label{fr-20-streaming-query-api-openai-compatible}

\textbf{Description:} The query flow shall be exposed as an HTTP service
that:

\begin{itemize}
\tightlist
\item
  Accepts chat-style requests in an \textbf{OpenAI-compatible} request
  format
\item
  Streams the generated response token-by-token via \textbf{Server-Sent
  Events (SSE)}
\item
  Returns source-attribution metadata alongside the streamed content
\end{itemize}

This contract is the integration surface for the Drupal front-end
(FR-16).

\textbf{Priority:} High\\
\textbf{Testability:} Send an OpenAI-format chat request to the
endpoint; verify an SSE stream is returned and that an OpenAI-compatible
client (or the Drupal AI provider) can consume it.

\subsubsection{FR-17: Vocabulary Mismatch
Bridging}\label{fr-17-vocabulary-mismatch-bridging}

\textbf{Description:} The retrieval layer shall surface documents
containing historical or domain-specific terminology when queried with
modern equivalent phrasing (e.g., a query for ``aviation'' retrieves
documents using ``flying machines'').

\textbf{Priority:} High\\
\textbf{Testability:} Define a vocabulary-mismatch test set; verify
recall@10 ≥ 0.8 for paired modern/historical query equivalents.

\subsubsection{FR-18: Document Removal from
Index}\label{fr-18-document-removal-from-index}

\textbf{Description:} The system shall support removing a document from
the retrieval index (deleting all its chunks from ChromaDB and the
metadata index) so that deprecated or restricted content is no longer
returned in responses.

\textbf{Priority:} Medium\\
\textbf{Testability:} Remove a document; verify no subsequent query
returns content from that document.

\subsubsection{FR-19: Ingestion Status
Visibility}\label{fr-19-ingestion-status-visibility}

\textbf{Description:} The system shall expose per-document ingestion
status to archivist/admin users: pending, processing, indexed, or
failed. Failed ingestions shall include a human-readable error message.

\textbf{Priority:} Low\\
\textbf{Testability:} Submit a document; observe status transitions from
pending → processing → indexed in the admin interface.

\subsection{3. Constraints}\label{constraints}

{\def\LTcaptype{none} % do not increment counter
\begin{longtable}[]{@{}
  >{\raggedright\arraybackslash}p{(\linewidth - 2\tabcolsep) * \real{0.5000}}
  >{\raggedright\arraybackslash}p{(\linewidth - 2\tabcolsep) * \real{0.5000}}@{}}
\toprule\noalign{}
\begin{minipage}[b]{\linewidth}\raggedright
ID
\end{minipage} & \begin{minipage}[b]{\linewidth}\raggedright
Constraint
\end{minipage} \\
\midrule\noalign{}
\endhead
\bottomrule\noalign{}
\endlastfoot
C-01 & OCR for scanned/image-based documents is out of scope; only
machine-readable text documents are ingested. \\
C-02 & LLM fine-tuning is out of scope; retrieval compensates for domain
knowledge gaps. \\
C-03 & All inference and document storage must remain on-premises; no
public cloud APIs for document content. \\
C-04 & Hardware target is a single GPU node with 24 GB VRAM (NVIDIA L4
or RTX 4090 equivalent). \\
C-05 & Python services use \texttt{pyproject.toml} for dependency/build
configuration and are served with Uvicorn (ASGI). \\
C-06 & The Drupal front-end uses the AI Chatblock module; integration
must work without a custom plugin, or scope expands (see FR-21). \\
C-07 & The LLM / inference runtime runs on the university GPU server,
accessed over SSH (e.g., an SSH tunnel to the model endpoint). The 24 GB
GPU node (C-04) refers to this server. \\
\end{longtable}
}

\subsection{4. Research Gaps (Future
Scope)}\label{research-gaps-future-scope}

The following are identified research gaps and are \textbf{not} part of
v1.0:

\begin{itemize}
\tightlist
\item
  Evaluation benchmarks specific to archival retrieval tasks
\item
  Handwritten and multimodal archival material processing
\item
  Provenance-aware answer generation
\item
  Privacy-preserving RAG
\item
  Explainable retrieval interfaces
\item
  Agentic archival research assistants (multi-step autonomous research)
\end{itemize}

\end{document}
